\documentclass[12pt, titlepage]{article}
%\usepackage[norsk]{babel}
\usepackage[utf8]{inputenc}
\usepackage{minted}
\usepackage[colorlinks = true,
linkcolor = blue,
urlcolor = blue]{hyperref}


\begin{document}

\title{Mappeoppgave - del 1}
\author{IIRA2001 - Høst 2023}
\date{}
\maketitle

%\section{Mappeoppgave 1}
Første del av mappen handler om generell programmering og tidsrekker i økonomiske fag.
Du skal legge ved tre små programmer om temaer fra faget så langt:
\begin{itemize}
  \item {Populasjonsvekst og den logistiske formelen}
  \item {Sparekalkulator}
  \item {Lånekalkulator}
\end{itemize}

Hvert av programmene trenger litt tekst som forklarer hva hensikten er eller hvilket problem det forsøker å løse,
samt hvordan det gjør dette, og eventuelt hva som skjer i ulike deler av programmet.

Dere kan for eksempel skrive denne teksten inn sammen med programmene i jupyter-notebook filer, eller
legge de med som et eget dokument om du har skrevet koden i en annen IDE som feks PyCharm.

I første del av mappen ser vi etter at dere kan : 
\begin{itemize}
  \item {lage egne funksjoner}
  \item {bruke variabler på en fornuftig måte}
  \item {behandle data ved bruk av lister eller dictionaries}
  \item {manipulere disse dataene med \mintinline{python}{for} eller \mintinline{python}{while} løkker}
  \item {styre programflyt med \mintinline{python}{if}/\mintinline{python}{else}}
  \item {plotte resultat med \mintinline{python}{pyplot}}
  \item {\textbf{Bruke python som verktøy i en faglig sammenheng}}
\end{itemize}

Dere \textbf{må} ikke bruke alle konseptene vi har sett på så langt, \textit{det aller viktigste er at dere kan 
bruke python til å løse et praktisk problem eller undersøke en relevant problemstilling.}

Det eneste vi forventer \textbf{skal} være med - er en fornuftig fremstilling av resultatene.
Det vil som oftest si et plot av en eller form, men om dette ikke passer, kunne det for eksempel
være en tabell eller en tekst 


I tillegg skal det være med et \textit{refleksjonsnotat}(se eget avsnitt).

\section*{Innhold}
\subsection*{Populasjonsvekst}
Oppgaver 1,3,5 og 9 har omhandlet populasjonsvekst og den logistiske formelen.
Som et minimum legger dere ved løsningen deres på en av disse oppgavene.

Dersom dere vil forbedre dette, kan dere ta i bruk programmet til å undersøke noe konkret.
Kan dere for eksempel si noe om populasjonsvekst på Sunnmøre eller i Norge?

Den logistiske likningen brukes også til å modellere hvordan ny teknologi tas i bruk (se \href{https://en.wikipedia.org/wiki/Diffusion_of_innovations}{Diffusion of innovation} og \href{https://no.wikipedia.org/wiki/Diffusjon_(kommunikasjon)}{Diffusjon (norsk wiki)})
Kanskje kan du si noe om hvor mange el-biler man kan forvente på veiene fremover, hvor mange flere brukere netflix kan håpe på å få, eller liknende.

\subsection*{Sparekalkulator}

Oppgaver 2,4,6,7 og 8 handlet om en sparekalkulator.
Som et minimum legger dere ved løsningen deres på en av disse oppgavene.

Som en forbedring, kan dere som foreslått i oppgaven også ta med inflasjon, formueskatt of skatt av rentene i sparekalkulatoren.
Kommer dere på mer nyttig funksjonalitet i et program som har med sparing å gjøre, kan dere implementere disse,
eller kanskje dere kan undersøke noen konkrete tilfeller som dere synes er interessante (Er du kanskje \href{https://www.imdb.com/title/tt0584425/}{Futurama} fan?)

\subsection*{Lånekalkulator}

Siste python-program i mappe 1 skal være en Lånekalkulator. Vi har ikke jobbet med lånekalkulatorer så langt, men programmet vil nok likne en del på sparekalkulatoren.
Her skal du lage et pythonprogram som gir deg en nedbetalingsplan og viser deg hvor mye et lån vil koste.

Du er fri til å designe programmet ditt, kanskje er den for boliglån slik som \href{https://www.nordea.no/privat/vare-produkter/lan-og-kreditt/boliglan/hva-koster-boliglanet.html#/}{her}, eller en serielånskalkulator slik som \href{https://www.smartepenger.no/kalkulatorer/978-serielanskalkulator}{her}

\subsection*{Forklarende tekst}

I tillegg til python-koden legger du ved en tekst til hver deloppgave, hvor du forklarer eventuell problemstilling, hva programmet ditt gjør og hvordan, og eventuelt drøfting av resultatene fra programmet.

\section*{Refleksjonsnotat}

\subsection*{Hva er et refleksjonsnotat?}

Et refleksjonsnotat er en tekst man ofte tilvirker når man er midtveis i eller ferdig med en prosess - ofte en form for praksis.
For eksempel kan en lærerstudent skrive et refleksjonsnotat hvor han reflekterer rundt en hendelse som skjedde da han var ute i praksis.
Hva skjedde, hvordan handlet du, hva lærte du og hva sier faglitteraturen. Refleksjonsnotater kan også gjelde for hele praksisperioden og omhandle faglig utvikling og veien videre.

\subsection*{Innholdet i refleksjonsnotat til første del av mappen}

Dere skal legge ved et refleksjonsnotat hvor dere reflekterer rundt læringsprosessen dere har vært gjennom, og rundt bruken av python i faget deres.
Typiske spørsmål dere reflekterer rundt vil kanskje være:
\begin{itemize}
   \item {Hva har jeg lært?}
   \item {Hva har vært vanskelig, og hvordan skal du håndtere det i videre virke?}
   \item {Med bakgrunn i oppgavene du har levert: hva vil komme til å være nyttig / mindre nyttig for deg?}
   \item {Hvilke utfordringer har \textit{du} med å bruke programmeringsverktøy in ditt fagområde fremover}
   \item {Hvordan tror du at du får bruk for verktøy som python videre}
   \item { ... }
\end{itemize}

Refleksjonsnotatet er altså en personlig tekst, og den skal vise faglig utvikling, eller refleksjon rundt faget deres sett sammenheng med programmering.
Det er bra om du knytter refleksjoner til konkrete eksempler fra kode-oppgavene dere har levert.

\section*{Innlevering og tilbakemelding}

Kode, kodetekst og refleksjonsnotat leveres på BlackBoard - Innleveringsfrist i (\color{blue}{slutten av uke 41?}\color{black}).
Deretter skal dere få mulighet til tilbakemelding på innleveringen slik at dere kan forbedre den inn mot \textit{endelig innlevering i inspera før jul} (Innleveringsfrist ikke satt ennå)

\end{document}
